% Generated draft: requires mathematical and source review.
\documentclass[11pt]{article}
\usepackage[margin=25mm]{geometry}
\usepackage{fontspec}
\setmainfont{Latin Modern Roman}
\usepackage{amsmath,amssymb,mathtools,unicode-math}
\setmathfont{Latin Modern Math}
\usepackage{enumitem,needspace,xurl,hyperref}
\setlength{\parindent}{0pt}
\setlength{\parskip}{5pt}
\setlength{\emergencystretch}{3em}
\newcommand{\N}{\mathbb{N}}
\newcommand{\Z}{\mathbb{Z}}
\newcommand{\Q}{\mathbb{Q}}
\newcommand{\R}{\mathbb{R}}
\newcommand{\C}{\mathbb{C}}
\newcommand{\F}{\mathbb{F}}
\newcommand{\A}{\mathbb{A}}
\newcommand{\PP}{\mathbb{P}}
\newcommand{\E}{\mathbb{E}}
\newcommand{\eps}{\varepsilon}
\newcommand{\dd}{\,\mathrm{d}}
\begin{document}
\section*{UnsolvedMath ID: 79}
\textbf{Source problem code:} catalogue entry
\subsection{Definitions and mathematical statement}
DARPA's Mathematical Challenge Ten (2007), ``Algorithmic Origami and Biology,'' asks for a stronger mathematical theory of isometric and rigid embedding relevant to protein folding. The supplied background links origami mathematics to folding at molecular scales and identifies drug design and protein engineering as possible applications.

For terminology only, an isometric embedding of metric spaces $(X,d_X)$ and $(Y,d_Y)$ is an injective map $f:X\to Y$ satisfying $d_Y(f(x),f(x'))=d_X(x,x')$ for every $x,x'\in X$. This definition does not select the intended embedding model: in folding problems, intrinsic surface isometry, prescribed framework distances, and rigid motions of faces impose different conditions. The source gives no spaces, dimensions, distance constraints, admissible deformations, collision conditions, or precise meaning of ``rigid embedding.''

The research task is to specify an appropriate embedding model and derive a protein-folding prediction or optimisation guarantee from it. No energy functional, input data, protein representation, target conformation, algorithmic guarantee, or quantified theorem is supplied. Thus the challenge is an application and theory-development agenda, not a single testable embedding statement. The registry open label and the partially-solved classification in the supplied triage dated 2026-08-17 are source classifications; that triage states that no unitary solved/open status or full solution/refutation is applicable as written, while noting separate literatures on rigidity, origami design, and computational protein modelling.
\subsection{Short English statement}
Develop isometric and rigid embedding theory that gives useful guarantees for protein folding, with possible applications to drug design and protein engineering. A precise problem still requires a folding model, constraints, and a prediction or optimisation objective.
\subsection{Sources}
\begin{itemize}
\item DARPA DSO, Broad Agency Announcement 07-68, Mathematical Challenge Ten (2007), supplied primary reference for the isometric/rigid-embedding agenda: \url{https://web.math.utk.edu/~vasili/refs/darpa07.MathChallenges.html}.
\item Local registry data/unsolvedmath/problems.json, permanent catalogue entry, source code catalogue entry, ``Algorithmic Origami and Biology,'' attributed to DARPA (2007); complete statement, background, and literature triage dated 2026-08-17. Dataset attribution: Ulam AI and the UnsolvedMath Contributors.
\end{itemize}
\end{document}
